%!TEX root = ../template.tex
%%%%%%%%%%%%%%%%%%%%%%%%%%%%%%%%%%%%%%%%%%%%%%%%%%%%%%%%%%%%%%%%%%%
%% chapter2.tex
%% NOVA thesis document file
%%
%% Chapter with state of the art
%%%%%%%%%%%%%%%%%%%%%%%%%%%%%%%%%%%%%%%%%%%%%%%%%%%%%%%%%%%%%%%%%%%

\typeout{NT FILE chapter2.tex}%

\chapter{State of the Art}
\label{cha:state_of_the_art}

This chapter surveys the existing landscape of \gls{GNSS} signal generation tools and related software interfaces, establishing the context in which the work of this internship was conducted. The goal is to identify the strengths and limitations of existing solutions and motivate the need for the developed system.

\section{Commercial GNSS Signal Simulators}

Commercial \gls{GNSS} simulators represent the industry standard for professional receiver testing and validation. These systems are purpose-built, highly capable, and widely used in aerospace, automotive, and defence industries.

The Spirent GSS7000 is one of the most prominent examples of a commercial \gls{GNSS} simulator. It supports simultaneous coherent signal generation across multiple constellations, including \gls{GPS}, \gls{GLONASS}, \gls{Galileo}, and \gls{BeiDou}, with up to 256 channels flexibly allocated across constellations \cite{spirentGSS7000}\cite{spirentGSS7000Federal}. The system is controlled via Spirent's SimGEN software, which provides comprehensive scenario definition, execution, and data management capabilities. While powerful, the GSS7000 is a closed, proprietary system with significant cost, and its control interface is tightly coupled to Spirent's own software ecosystem, offering limited flexibility for programmatic integration into custom workflows.

Similar offerings exist from other vendors such as Rohde \& Schwarz \cite{rohdeSchwarzSMBV100B} and Keysight Technologies \cite{keysightN7609B}, which provide high-fidelity \gls{GNSS} signal generation equipment. These systems share the same fundamental limitations: they are expensive, closed-source, and not designed for flexible programmatic control by third-party software. They are typically operated through vendor-specific graphical interfaces, making automated testing pipelines difficult to implement without additional integration effort.

\section{Open-Source GNSS Signal Generation Tools}

\textbf{REVIEW}

In contrast to commercial solutions, several open-source tools have emerged that leverage \gls{SDR} hardware to generate \gls{GNSS} signals at a fraction of the cost.

\textbf{GPS-SDR-SIM} \cite{gpsSdrSim} is a widely adopted open-source GPS signal simulator. It generates GPS L1 C/A baseband signal streams as digitised I/Q samples, which can then be transmitted via \gls{SDR} hardware such as HackRF or ADALM-Pluto. The tool accepts a user-defined trajectory specified as a CSV file containing Earth-Centred Earth-Fixed (ECEF) positions, or as an NMEA GGA stream, and uses a GPS broadcast ephemeris file to compute pseudoranges and Doppler shifts for satellites in view. Despite its usefulness, GPS-SDR-SIM has notable limitations: it is restricted to GPS L1 C/A signals, operates primarily in an offline batch mode generating pre-computed I/Q sample files, and provides no programmatic API for runtime parameter control.

\textbf{Galileo-SDR-SIM} \cite{galileoSdrSim} extends the GPS-SDR-SIM concept to the Galileo E1B/C signal, generating Galileo baseband signal streams for transmission via \gls{SDR} platforms such as the USRP B210. While it demonstrates the feasibility of open-source multi-constellation simulation, it remains a standalone command-line tool with limited runtime configurability and no structured interface for integration into larger software systems.

\textbf{GNSS-SDR} \cite{gnssSdr} is an open-source \gls{GNSS} receiver framework rather than a signal generator. It is designed for signal reception and processing, not generation, and therefore does not directly address the use case of the present work. Its architecture is nonetheless instructive, demonstrating how modular \gls{SDR}-based \gls{GNSS} software can be structured.

A common limitation across all open-source tools is their focus on specific \gls{SDR} hardware platforms such as HackRF, ADALM-Pluto, or USRP devices. None provide support for professional-grade hardware such as DekTec \gls{SDR} cards, nor do they expose a structured, remotely accessible API for programmatic control.

\section{Hardware Vendor APIs}

DekTec produces a range of professional \gls{SDR} hardware, including the DTA-2115 and DTA-2116 PCIe cards used in this internship. These cards are capable of high-quality, high-bandwidth RF signal generation and are well suited for professional signal simulation applications.

DekTec provides the DTAPI, a C++ library for interfacing with their hardware. While functional and comprehensive in its coverage of hardware capabilities, the DTAPI is a low-level interface that exposes hardware primitives directly. Configuring a device requires a detailed understanding of internal hardware parameters, correct sequencing of API calls, and careful management of device state. There is no existing higher-level abstraction layer for \gls{GNSS} use cases, nor any remote control interface built on top of the DTAPI. Developers must therefore deal with low-level concerns such as FIFO management, DMA buffer sizing, and modulation flag composition directly, which increases both development time and the risk of subtle bugs.

\section{Summary and Identified Gap}

Table~\ref{tab:sota_comparison} summarises the key characteristics of the surveyed solutions.

\begin{table}[h]
\centering
\caption{Comparison of existing GNSS signal generation solutions.}
\label{tab:sota_comparison}
\begin{tabular}{|l|c|c|c|c|}
    \hline
    \textbf{Solution} & \textbf{Multi-constellation} & \textbf{Programmatic API} & \textbf{Open Source} & \textbf{DekTec Support} \\
    \hline
    Spirent GSS7000   & \checkmark & Limited  & \texttimes & \texttimes \\
    GPS-SDR-SIM       & \texttimes & \texttimes & \checkmark & \texttimes \\
    Galileo-SDR-SIM   & \texttimes & \texttimes & \checkmark & \texttimes \\
    GNSS-SDR          & \checkmark & \texttimes & \checkmark & \texttimes \\
    DekTec DTAPI      & \checkmark & Low-level only & \texttimes & \checkmark \\
    \hline
    \textbf{This work} & \checkmark & \checkmark & \texttimes & \checkmark \\
    \hline
\end{tabular}
\end{table}

The survey reveals a clear gap in the existing landscape: no available solution combines support for professional DekTec hardware with a clean, high-level programmatic interface suitable for integration into automated testing workflows. Commercial simulators are closed and costly; open-source tools target consumer \gls{SDR} hardware and lack runtime configurability; and the DekTec DTAPI, while powerful, requires significant low-level expertise to use correctly.

The work developed during this internship directly addresses this gap, providing a C++ abstraction layer over the DTAPI, exposed through a \gls{gRPC} interface, enabling programmatic and remotely accessible control of \gls{GNSS} signal generation on professional-grade hardware.

\section{Conclusion}

This chapter introduce the state of the art in \gls{GNSS} signal generation, covering commercial simulators, open-source tools, and hardware vendor APIs. The analysis identified a gap between low-level hardware access and the need for a flexible, programmatically controllable signal generation interface. This gap motivates the development described in the following chapters.
